\subsection{Hensel 引理}

在拓扑环 A 中，若元素 x 满足 0 是序列 \((x^{n})_{n\geqslant0}\) 的极限，则称 x 是拓扑幂零的。若 A 是具有线性拓扑的交换环，则称 x 拓扑幂零，是指对于 A 的每个开理想 \(\mathfrak{S}\)，x 在 A/\(\mathfrak{S}\) 中的典范像是该环的幂零元。若 \(r_{\mathfrak{S}}\) 是 A/\(\mathfrak{S}\) 的幂零根，则显然 \((r_{\mathfrak{S}})\) 是子集的逆系统，而 A 的拓扑幂零元集 t 是 \(r=\lim_{\leftarrow \mathfrak{S}}r_{\mathfrak{S}}\) 在

典范同态 \(\mathbf{A} \to \lim_{\leftarrow \mathfrak{S}} \mathbf{A} / \mathfrak{S}\) 下的逆像；因此它是 \(\mathbf{A}\) 的闭理想。若进一步 \(\mathbf{A}\) 是 Hausdorff 且完备的，则此理想包含于 \(\mathbf{A}\) 的 Jacobson 根中，并且元素 \(x \in \mathbf{A}\) 可逆的充要条件是其模 \(t\) 的类在 \(\mathbf{A} / t\) 中可逆（第 2 节，第 13 号，引理 3）。

注意，若 A 是环且 m 是 A 的双边理想，则 m 中的元素关于 m-进拓扑都是拓扑幂零的。

定理 1（Hensel）。设 A 为完备的 Hausdorff 线性拓扑交换环。设 m 是 A 的闭理想，其元素都是拓扑幂零的。令 B = A/m 为商拓扑环，\(\varphi: A \to B\) 为典范映射。令 R 为 A\{X\} 中的限制形式幂级数，\(\overline{\mathrm{P}}\) 为 B{[}X{]} 中的首一多项式，\(\overline{\mathrm{Q}}\) 为 B\{X\} 中的限制形式幂级数。假设 \(\overline{\varphi}(\mathrm{R}) = \overline{\mathrm{P}}\) 。\(\overline{\mathrm{Q}}\)，并且 \(\overline{\mathrm{P}}\) 与 \(\overline{\mathrm{Q}}\) 在 B\{X\} 中强互素。则存在唯一的有序对 (P, Q)，其中 \(\mathbf{P} \in \mathbf{A}[\mathbf{X}]\) 是首一多项式，\(\mathbf{Q} \in \mathbf{A}\{\mathbf{X}\}\) 是限制形式幂级数，并满足

\[
\mathrm{R} = \mathrm{P}. \mathrm{Q}, \quad \overline {{{\varphi}}} (\mathrm{P}) = \overline {{{\mathrm{P}}}}, \quad \overline {{{\varphi}}} (\mathrm{Q}) = \overline {{{\mathrm{Q}}}}.\tag{4}
\]

此外，P 与 Q 在 A\{X\} 中强互素；若 R 是多项式，则 Q 也是多项式。

证明分为四步。在前三步中假设 A 是离散的，此时 R 和 \(\overline{Q}\) 是多项式。

\textbf{(1) \(\mathfrak{m}^2 = 0\)}\par

设 S、T 是 A{[}X{]} 中的两个多项式，S 是首一的，且 \(\overline{\varphi}(S) = \overline{P}\)、\(\overline{\varphi}(T) = \overline{Q}\)；第 1 号命题 2 表明 S、T 强互素；所以（第 1 号，命题 1）存在唯一的 A{[}X{]} 中的多项式有序对 (S', T')，使得

\[
\mathrm{R} - \mathrm{ST} = \mathrm{ST} ^ {\prime} + \mathrm{TS} ^ {\prime} \quad \text { and } \quad \deg (\mathrm{S} ^ {\prime}) <   \deg (\mathrm{S}) = \deg (\overline {{\mathrm{P}}}).\tag{5}
\]

多项式 \(P = S + TS'\)、\(Q = T + T'\) 就是问题的解；事实上

\[
\overline {{{\mathbf {P}}}} \cdot \overline {{{\varphi}}} (\mathbf {T} ^ {\prime}) + \overline {{{\mathbf {Q}}}} \cdot \overline {{{\varphi}}} (\mathbf {S} ^ {\prime}) = \overline {{{\varphi}}} (\mathbf {S T} ^ {\prime} + \mathbf {T S} ^ {\prime}) = \overline {{{\varphi}}} (\mathbf {R} - \mathbf {S T}) = 0.\tag{6}
\]

由于 \(\overline{\mathbf{P}}\) 是首一的，\(\overline{\mathbf{P}}\) 与 \(\overline{\mathbf{Q}}\) 强互素，且

\[
\deg (\overline {{{\varphi}}} (S ^ {\prime})) <   \deg (\overline {{{\mathbf {P}}}}),
\]

第 1 号命题 1 表明 \(\bar{\varphi}(\mathbf{S}') = \bar{\varphi}(\mathbf{T}') = 0\)，换言之，\(\mathbf{S}'\) 和 \(\mathbf{T}'\) 的系数属于 \(\mathfrak{m}\)，而关系 \(\mathfrak{m}^2 = 0\) 给出

\[
\mathrm{PQ} = \mathrm{ST} + \mathrm{ST} ^ {\prime} + \mathrm{TS} ^ {\prime} = \mathrm{R},
\]

这满足关系（4）。由于 \(\overline{\varphi}(\mathbf{P}) = \overline{\mathbf{P}}\) 且 \(\overline{\varphi}(\mathbf{Q}) = \overline{\mathbf{Q}}\)，\(\mathbf{P}\) 与 \(\mathbf{Q}\) 强互素（第 1 号，命题 2）；最后，若 \(\mathbf{P}_1\) 和 \(\mathbf{Q}_1\) 是 A{[}X{]} 中满足（4）的另外两个多项式，且 \(\mathbf{P}_1\) 是首一的，则必有 \(S_1' = P_1 - S\)、\(T_1' = Q_1 - T\)、\(\deg(S_1') < \deg(S)\) 以及 \(R - ST = ST_1' + TS_1'\)，因为 \(S_1'\) 和 \(T_1'\) 的系数属于 \(m\)；但命题 1 随即证明 \(S' = S_1'\) 且 \(T' = T_1'\)，这就证明了有序对 \((P, Q)\) 的唯一性。

\textbf{(2) \(\mathfrak{m}\) 是幂零的}\par

令 \(n\) 为满足 \(m^n = 0\) 的最小整数，并对 \(n > 2\) 作归纳论证；定理在 \(n = 2\) 时已经证明。令 \(A' = A / m^{n-1}\)、\(m' = m / m^{n-1}\)；由于 \(m'^{n-1} = 0\)，存在唯一的有序对 \((P', Q')\)，其元素是 \(A'[X]\) 中的多项式，使得 \(P'\) 是首一的，\(R' = P'Q'\)，\(\bar{\psi}(P') = \overline{P}\) 且 \(\bar{\psi}(Q') = \overline{Q}\)，其中 \(\psi\) 表示典范同态 \(A' \to A'/m' = B\)，\(\theta\) 表示典范同态 \(A \to A'\)，且 \(R' = \bar{\theta}(R)\)。另一方面，由于 \((\mathfrak{m}^{n-1})^{2}=0\)，存在唯一的有序对 \((\mathrm{P},\mathrm{Q})\)，其元素是 A{[}X{]} 中的多项式，使得 P 是首一的且 \(R=PQ,\bar{\theta}(P)=\overline{P},\bar{\theta}(Q)=\overline{Q}\)；由于 \(\phi=\psi\circ\theta\)，这表明满足（4）的 P、Q 存在且唯一；此外，根据归纳假设，\(P'\) 与 \(Q'\) 强互素，因而 P 与 Q 也强互素。

\textbf{(3) A 是离散的}\par

注意，在此情形 m 不再必然是幂零的，但根据假设它总是幂理想。设 \(P_{0}\)、\(Q_{0}\) 是 A{[}X{]} 中的两个多项式，满足 \(\overline{\varphi}(P_{0}) = \overline{P}\)、\(\overline{\varphi}(Q_{0}) = \overline{Q}\)，且 \(P_{0}\) 是首一的。考虑 A 中由 \(R - P_{0}Q_{0}\) 的系数生成的理想 n；它是有限生成的且包含于 m，因而是幂零的（第二章，第 2 节，第 6 号，命题 15）；按定义，若 \(\psi: A \to A/n\) 是典范映射，则 \(\overline{\psi}(R) = \overline{\psi}(P_{0})\overline{\psi}(Q_{0})\)。此外，\(\overline{\psi}(P_{0})\) 与 \(\overline{\psi}(Q_{0})\) 强互素；这是由 P、Q 的假设以及将第 1 号命题 2 应用于典范同态 \(A/n \to A/m\) 得出的。根据情形（2），存在有序对 \((P, Q)\)，其元素是 A{[}X{]} 中的多项式，使得 P 是首一的且关系式（4）成立。P 与 Q 强互素这一事实同样蕴含 P 与 Q 在 A{[}X{]} 中强互素，因为 m 包含于 A 的 Jacobson 根中（第 1 号，命题 2）。最后，设 \(P_{1}\)、\(Q_{1}\) 是两个多项式，属于 A{[}X{]}，满足（4）且 \(P_{1}\) 是首一的；令 \(n_{1}\) 为 A 中由 \(P - P_{1}\) 的系数和 \(Q - Q_{1}\) 的系数生成的有限生成理想；由于 \(n_{1}\) 包含于 m，它是幂零的；若 \(\psi_{1}: A \to A/n_{1}\) 是典范映射，则 \(\overline{\psi}_{1}(P) = \overline{\psi}_{1}(P_{1})\) 且

\[
\overline {{{{\psi}}}} _ {1} (\mathbf {Q}) = \overline {{{{\psi}}}} _ {1} (\mathbf {Q} _ {1});
\]

因此，情形（2）的唯一性蕴含 \(\mathrm{P} = \mathrm{P}_1, \mathrm{Q} = \mathrm{Q}_1\)。

\textbf{(4) 一般情形}\par

令 \(\mathcal{B}\) 为 A 中由 A 的理想组成的 0 的邻域基本系。对于所有 \(\mathfrak{S} \in \mathcal{B}\)，令 \(f_{\mathfrak{S}}\) 为典范映射 \(\mathrm{A} \to \mathrm{A} / \mathfrak{S}\)，令 \(\phi_{\mathfrak{S}}\) 为典范映射

\[
\mathrm{A} / \mathfrak{S} \mapsto (\mathrm{A} / \mathfrak{S}) / ((\mathfrak {m} + \mathfrak{S}) / \mathfrak{S}) = \mathrm{A} / (\mathfrak {m} + \mathfrak{S}),
\]

\(g_{\mathfrak{S}}\) 为典范映射 \(\mathrm{B} = \mathrm{A} / \mathfrak{m}\to \mathrm{A} /(\mathfrak{m} + \mathfrak{S})\)，并记 \(\mathrm{R}_{\mathfrak{S}} = \bar{f}_{\mathfrak{S}}(\mathrm{R})\)，\(\overline{\mathrm{P}}_{\mathfrak{S}} = \bar{g}_{\mathfrak{S}}(\overline{\mathrm{P}}),\overline{\mathrm{Q}}_{\mathfrak{S}} = \bar{g}_{\mathfrak{S}}(\overline{\mathrm{Q}})\)。由于每个环 \(\mathrm{A} / \mathfrak{S}\) 都是离散的，故可将情形 (3) 应用于它，于是可知存在唯一的有序对 \((\mathrm{P}_{\mathfrak{S}},\mathrm{Q}_{\mathfrak{S}})\)，其元素是 \((\mathrm{A} / \mathfrak{S})[\mathrm{X}]\) 中的多项式，使得 \(\mathrm{P}_{\mathfrak{S}}\) 是首一的，且 \(\mathrm{R}_{\mathfrak{S}} = \mathrm{P}_{\mathfrak{S}}\mathrm{Q}_{\mathfrak{S}},\overline{\varphi}_{\mathfrak{S}}(\mathrm{P}_{\mathfrak{S}}) = \overline{\mathrm{P}}_{\mathfrak{S}},\overline{\varphi}_{\mathfrak{S}}(\mathrm{Q}_{\mathfrak{S}}) = \overline{\mathrm{Q}}_{\mathfrak{S}}\)。这个有序对的唯一性意味着：若 \(\mathfrak{S}'\subset \mathfrak{S},\mathfrak{S}'\in \mathcal{B}\)，且 \(f_{\mathfrak{S}\mathfrak{S}'}:\mathrm{A} / \mathfrak{S}'\rightarrow \mathrm{A} / \mathfrak{S}\) 是典范映射，则 \(\mathrm{P}_{\mathfrak{S}} = \bar{f}_{\mathfrak{S}\mathfrak{S}'}(\mathrm{P}_{\mathfrak{S}'})\)，\(\mathrm{Q}_{\mathfrak{S}} = \bar{f}_{\mathfrak{S}\mathfrak{S}'}(\mathrm{Q}_{\mathfrak{S}'})\)。于是由 \(\mathrm{A}\{\mathrm{X}\}\) 与 \(\lim (\mathrm{A} /\mathfrak{S})[\mathrm{X}]\) 的典范同一

\[
\varprojlim_ {\mathfrak {S}} (\mathrm{A} / \mathfrak {S}) [ X ]
\]

(第 2 号的命题 3) 可知，存在 \(\mathbf{P} \in \mathbf{A}\{\mathbf{X}\}\) 和 \(\mathbf{Q} \in \mathbf{A}\{\mathbf{X}\}\)，使得 \(\mathbf{R} = \mathbf{PQ}\)，并且 \(\bar{f}_{\mathfrak{S}}(\mathbf{P}) = \mathbf{P}_{\mathfrak{S}}, \bar{f}_{\mathfrak{S}}(\mathbf{Q}) = \mathbf{Q}_{\mathfrak{S}}\) 对所有 \(\mathfrak{S} \in \mathcal{B}\) 都成立。此外，

\[
\bar {g} _ {\mathfrak {S}} (\overline {{{\mathbf {P}}}} - \overline {{{\varphi}}} (\mathbf {P})) = 0, \quad \bar {g} _ {\mathfrak {S}} (\overline {{{\mathbf {Q}}}} - \overline {{{\varphi}}} (\mathbf {Q})) = 0
\]

对所有 \(\mathfrak{S} \in \mathcal{B}\) 都成立，这意味着：对于所有 \(\mathfrak{S} \in \mathcal{B}\)，\(\overline{\mathbf{P}} - \overline{\varphi}(\mathbf{P})\) 和 \(\overline{\mathbf{Q}} - \overline{\varphi}(\mathbf{Q})\) 的所有系数都属于 \((\mathfrak{m} + \mathfrak{S})/\mathfrak{m}\)。但是，由于 \(\mathfrak{m}\) 在 \(\mathbf{A}, \bigcap_{\mathfrak{S}} (\mathfrak{m} + \mathfrak{S}) = \mathfrak{m}\) 中是闭的，从而 \(\overline{\mathbf{P}} = \overline{\varphi}(\mathbf{P}), \overline{\mathbf{Q}} = \overline{\varphi}(\mathbf{Q})\)，于是 \(\mathbf{P}\) 和 \(\mathbf{Q}\) 显然满足 (4)；此外，由于 \(\mathbf{P}_{\mathfrak{S}}\) 是首一的且次数相同，受限形式幂级数 \(\mathbf{P}\) 是首一多项式。若 \((\mathbf{P}', \mathbf{Q}')\) 是另一个满足 (4) 的有序对，且 \(\mathbf{P}'\) 是首一多项式，则可推出 \(\mathrm{R}_{\mathfrak{S}} = f_{\mathfrak{S}}(\mathrm{P}') f_{\mathfrak{S}}(\mathrm{Q}')\)，\(\overline{\varphi}_{\mathfrak{S}}(f_{\mathfrak{S}}(\mathrm{P}')) = \overline{\mathrm{P}}_{\mathfrak{S}}\) 以及 \(\overline{\varphi}_{\mathfrak{S}}(f_{\mathfrak{S}}(\mathrm{Q}')) = \overline{\mathrm{Q}}_{\mathfrak{S}}\)，由情形 (3) 中的唯一性可得 \(f_{\mathfrak{S}}(\mathrm{P}') = \mathrm{P}_{\mathfrak{S}}\)，\(f_{\mathfrak{S}}(\mathrm{Q}') = \mathrm{Q}_{\mathfrak{S}}\) 对所有 \(\mathfrak{S} \in \mathcal{B}\) 都成立，从而 \(\mathrm{P} = \mathrm{P}'\) 且 \(\mathrm{Q} = \mathrm{Q}'\)。最后证明 \(\mathbf{P}\) 和 \(\mathbf{Q}\) 强互素；根据情形 (3) 以及第 1 号的命题 1，对所有 \(\mathfrak{S} \in \mathcal{B}\)，存在唯一的有序对 \((\mathrm{S}_{\mathfrak{S}}, \mathrm{T}_{\mathfrak{S}})\)，其元素是 \((\mathrm{A}/\mathfrak{S})[\mathrm{X}]\) 中的多项式，使得

\[
1 = \mathrm{P} _ {\mathfrak {S}} \mathrm{S} _ {\mathfrak {S}} + \mathrm{Q} _ {\mathfrak {S}} \mathrm{T} _ {\mathfrak {S}} \quad \text { and } \quad \deg (\mathrm{T} _ {\mathfrak {S}}) <   \deg (\mathrm{P} _ {\mathfrak {S}}) = \deg (\bar {\mathrm{P}}).\tag{7}
\]

这个有序对的唯一性立即表明，对于 \(\mathfrak{S}' \in \mathcal{B}\)，\(\mathfrak{S}' \subset \mathfrak{S}\)，有 \(S_{\mathfrak{S}} = \bar{f}_{\mathfrak{S}\mathfrak{S}'}(S_{\mathfrak{S}'})\)，\(T_{\mathfrak{S}} = \bar{f}_{\mathfrak{S}\mathfrak{S}'}(T_{\mathfrak{S}'})\)；考虑第 2 号的命题 3，可得存在两个受限形式幂级数 \(S\)、\(T\)，属于 \(A\{X\}\)，使得 \(S_{\mathfrak{S}} = \bar{f}_{\mathfrak{S}}(S)\)，\(T_{\mathfrak{S}} = \bar{f}_{\mathfrak{S}}(T)\)，并且 \(1 = PS + QT\)。

最后还需验证：若 R 是多项式，则 Q 也是多项式。现在，按构造 \(Q_{\mathfrak{S}}\) 是多项式，并且由于 \(P_{\mathfrak{S}}\) 是首一的，关系 \(R_{\mathfrak{S}} = P_{\mathfrak{S}}Q_{\mathfrak{S}}\) 蕴含

\[
\deg (Q _ {\mathfrak {S}}) \leqslant \deg (R _ {\mathfrak {S}}) \leqslant \deg (R)
\]

对所有 \(\mathfrak{S} \in \mathcal{B}\) 都成立；由 \(\mathbf{Q}\) 的定义，立即得到所需结论。
% semantic-fragments: legacy-input-seam
\relax{}
